Un núvol elèctricament carregat està situat a 4,7~km d'altura sobre el terra. La diferència de potencial entre la base del núvol i el terra és de $2,3\times 10^{6}\ \mathrm{V}$. Suposem que el camp elèctric en aquesta regió és uniforme i que la càrrega elèctrica del núvol és positiva.

Una gota d'aigua que es troba entre el núvol i el terra té una massa d'1,3~mg i una càrrega de valor $Q$. En un moment donat, la gota ascendeix cap al núvol amb una velocitat constant de $2\ \mathrm{m\,s^{-1}}$ (sense tenir en compte els corrents d'aire ni el fregament).

\begin{apartat}{1}
Dibuixeu un esquema de la situació descrita pel problema i representeu-hi les càrregues elèctriques implicades i els camps vectorials (gravitatori i elèctric). Calculeu la intensitat del camp elèctric que hi ha entre el núvol i el terra, i indiqueu-ne el mòdul, la direcció i el sentit.
\end{apartat}

\begin{apartat}{1}
Calculeu el valor de la càrrega $Q$ (en nC) i expliqueu raonadament quin signe hauria de tenir.
\end{apartat}

\blocetiquetat{Dada:}{$g = 9,81\ \mathrm{m\,s^{-2}}$.}
